\chapter{Further Reading}
\label{app:reading}

\section{The primary sources}

\begin{description}
\item[\emph{Smalltalk-80: The Language and its Implementation},] Goldberg and
  Robson, 1983. The Blue Book. Part One is still the best introduction to the
  language, and Part Four is the specification this virtual machine is checked
  against. Chapter 30 is \ct|om_bb.c|.
\item[\emph{Smalltalk-80: The Interactive Programming Environment},] Goldberg,
  1984. The Orange Book, and the manual for the interface in
  Part~IV of this tutorial. If the Browser or the Debugger puzzles you, this is
  where the answer is.
\item[\emph{Multiprocessor Smalltalk},] Pallas and Ungar, PLDI 1988. Where the
  \emph{serialize / replicate / reorganize} taxonomy comes from, and still the
  best account of what goes wrong when a Smalltalk library meets real
  parallelism.
\end{description}

\section{This project's own documentation}

\begin{center}
\small
\begin{tabular}{@{}lp{0.58\textwidth}@{}}
\toprule
\ctp{doc/CONCURRENCY.md} & the contract, normative. The source for
  Chapter~\ref{ch:contract}\\
\ctp{doc/SCALING.md} & the benchmark, and what each kernel is actually
  measuring\\
\ctp{doc/Display.md} & the window, the face, antialiasing, and what the
  interface expects of you\\
\ctp{doc/LanguageExtensions.md} & every post-1983 syntax and where it stands\\
\ctp{doc/MultiThreading.md} & the threading implementation\\
\ctp{doc/PLAN-TO-PHARO.md} & where the project is going, sized honestly\\
\ctp{doc/LICENSING.md} & what may be redistributed, and what may not\\
\ctp{doc/PORTABILITY.md} & the platform layer\\
\bottomrule
\end{tabular}
\end{center}

\section{The best documentation is the image}

This is not a platitude. The 1983 library is 226 classes of small, readable,
commented Smalltalk written by people who were inventing the ideas as they
went, and the Browser is built to read it.

Three habits worth forming:

\begin{description}
\item[\ctp{implementors}] on any selector you do not understand. Reading three
  implementations of \ct{do:} teaches more about the collection hierarchy than
  any diagram.
\item[\ctp{senders}] on any method you are about to change. It is the whole of
  impact analysis, and it is instant.
\item[Read a class comment before its methods.] The 1983 comments say why the
  class exists, which is the part that is hard to reconstruct.
\end{description}

Start with \ct{Collection}, \ct{SequenceableCollection} and \ct{Stream}. Then
\ct{Object}, which is smaller than you expect. Then, when you want to see how
the interface fits together, \ct{Controller} and \ct{View} in
\ct{Interface-Framework}.

\section{The source tree}

\begin{center}
\small
\begin{tabular}{@{}lp{0.58\textwidth}@{}}
\toprule
\ctp{lib/Kernel-Exceptions/} & the exception system, in readable Smalltalk\\
\ctp{lib/Concurrency/} & \ctp{Mutex}, \ctp{Monitor}, \ctp{SharedQueue},
  \ctp{Promise}. The class comments explain the design decisions\\
\ctp{lib/Probe/} & small classes that exist to exercise one feature each ---
  traits, pragmas, unwinding, initialisation\\
\ctp{profiles/} & each profile's comment explains what it composes and why
  anything is excluded. They are worth reading as a set\\
\ctp{tests/unit/} & what the system checks about itself\\
\bottomrule
\end{tabular}
\end{center}

\section{Where to ask}

\begin{center}
\url{https://github.com/blakemcbride/Smalltalk-2026}
\end{center}
